\section*{Bab \ref{ch:b1:lewis-resonance-vsepr} --- Struktur Lewis, Resonansi, dan VSEPR}

\begin{solution}{exo:b1:lewis-resonance-vsepr:1}
\ce{H2O}: dua ikatan \ce{O-H}, dua \omterm{def:b1:lewis-resonance-vsepr:lewis}{pasangan elektron bebas} pada O.
\ce{NH3}: tiga ikatan \ce{N-H}, satu \omterm{def:b1:lewis-resonance-vsepr:lewis}{pasangan bebas} pada N. \ce{CO2}:
\ce{O=C=O}, dua \omterm{def:b1:lewis-resonance-vsepr:lewis}{pasangan bebas} pada setiap O. \ce{HCN}: \ce{H-C#N}, satu
\omterm{def:b1:lewis-resonance-vsepr:lewis}{pasangan bebas} pada N. \ce{CH2O}: karbon terikat pada dua H dan berikatan
rangkap dua dengan O, dua \omterm{def:b1:lewis-resonance-vsepr:lewis}{pasangan bebas} pada O. Semua \omterm{def:b1:lewis-resonance-vsepr:formal-charge}{muatan formal}
kelima molekul ini nol. \ce{NH4+}: empat ikatan \ce{N-H}, tanpa \omterm{def:b1:lewis-resonance-vsepr:lewis}{pasangan
bebas}; nitrogen $5 - 0 - 4 = +1$, yaitu muatan ion itu.
\end{solution}

\begin{solution}{exo:b1:lewis-resonance-vsepr:2}
\ce{H2S}: \ce{AX2E2}, bengkok, di bawah $109.5^\circ$ (terukur $92^\circ$).
\ce{PCl3}: \ce{AX3E}, piramida segitiga, sekitar $100^\circ$.
\ce{BF3}: \ce{AX3}, segitiga datar, $120^\circ$. \ce{SiH4}: \ce{AX4},
tetrahedral, $109.5^\circ$. \ce{CO3^{2-}}: \ce{AX3} (ikatan rangkap
dihitung sekali), segitiga datar, $120^\circ$.
% ledger: geo:H2S.aHSH, geo:PCl3.aClPCl
\end{solution}

\begin{solution}{exo:b1:lewis-resonance-vsepr:3}
Nonpolar: \ce{CCl4} (\ce{AX4} tetrahedral), \ce{BF3} (\ce{AX3}),
\ce{CO2} (\ce{AX2} linear): dipol-dipol ikatannya saling meniadakan.
Polar: \ce{CH2Cl2} (dua jenis ikatan yang berbeda pada sebuah
tetrahedron), \ce{NH3} (piramida), \ce{SO2} (bengkok).
\end{solution}

\begin{solution}{exo:b1:lewis-resonance-vsepr:4}
\ce{H+} adalah \omterm{def:b1:lewis-resonance-vsepr:lewis-acid}{asam Lewis} dan \ce{H2O} basanya: ikatan \ce{O-H} baru pada
\ce{H3O+} bersifat datif ketika terbentuk (sesudahnya ketiga ikatan
\ce{O-H} identik). \ce{Cu^{2+}} adalah asamnya dan setiap \ce{NH3} sebuah
basa: keempat ikatan \ce{Cu-N} bersifat datif.
\end{solution}

\begin{solution}{exo:b1:lewis-resonance-vsepr:5}
\chemfig{CH_3-C(=[:60]O)-[:-60]\chemabove{O}{\scriptstyle\ominus}}
$\leftrightarrow$
\chemfig{CH_3-C(-[:60]\chemabove{O}{\scriptstyle\ominus})=[:-60]O}.
Kedua struktur itu ekuivalen: kedua ikatan \ce{C-O} berorde 1.5 dan sama
panjang, di antara ikatan rangkap dua dan ikatan tunggal. Pada asam
metanoat kedua oksigennya berbeda (satu membawa H), struktur-strukturnya
tidak ekuivalen, dan ikatan-ikatannya tetap berbeda panjang, 120.2
(\ce{C=O}) dan \qty{134.3}{pm} (\ce{C-OH}).
\end{solution}

\begin{solution}{exo:b1:lewis-resonance-vsepr:6}
$N = 6 + 4 \times 6 + 2 = 32$ elektron. Dengan empat ikatan tunggal dan
tiga \omterm{def:b1:lewis-resonance-vsepr:lewis}{pasangan bebas} pada setiap oksigen: belerang $6 - 0 - 4 = +2$, setiap
oksigen $6 - 6 - 1 = -1$; jumlahnya $+2 - 4 = -2$. Bentuk \ce{AX4},
tetrahedral. Dengan dua ikatan \ce{S=O}, memilih dua dari keempat oksigen
yang berikatan rangkap dua menghasilkan $\binom{4}{2} = 6$ struktur.
\end{solution}

\begin{solution}{exo:b1:lewis-resonance-vsepr:7}
\ce{SF4}: \ce{AX4E}, jungkat-jungkit. \ce{ClF3}: \ce{AX3E2}, bentuk T.
\ce{XeF2}: \ce{AX2E3}, linear. \ce{XeF4}: \ce{AX4E2}, segi empat datar.
Dalam bipiramida segitiga, posisi aksial mempunyai tiga tetangga pada
$90^\circ$, posisi ekuatorial hanya dua: pasangan-pasangan bebas yang
besar menempati tempat yang tolakan $90^\circ$-nya paling sedikit.
\end{solution}

\begin{solution}{exo:b1:lewis-resonance-vsepr:8}
Belerang dan fosfor lebih besar dan kurang elektronegatif daripada
oksigen dan nitrogen: pasangan-pasangan elektron ikatan terletak lebih
jauh dari atom pusat dan tertarik ke arah hidrogen, sehingga kurang
saling tolak, dan pasangan-pasangan bebas menekan ikatan-ikatan itu ke
arah $90^\circ$.
\end{solution}

\begin{solution}{exo:b1:lewis-resonance-vsepr:9}
\ce{CH3-C#N}: karbon \ce{CH3} sp$^3$, karbon nitril sp, nitrogen sp;
5 ikatan $\sigma$ (tiga \ce{C-H}, \ce{C-C}, satu dalam \ce{C#N}) dan
2 ikatan $\pi$. \ce{CH2=CH-CH3}: kedua karbon alkena sp$^2$, karbon
metil sp$^3$; 8 ikatan $\sigma$ (enam \ce{C-H}, dua \ce{C-C}) dan
1 ikatan $\pi$.
\end{solution}

\begin{solution}{exo:b1:lewis-resonance-vsepr:10}
Kedua molekul berbentuk piramida dengan \omterm{def:b1:lewis-resonance-vsepr:lewis}{sepasang elektron bebas} pada
nitrogen. \omterm{def:b1:lewis-resonance-vsepr:lewis}{Pasangan bebas} itu memberi sumbangan yang mengarah (dari
negatif ke positif) dari \omterm{def:b1:lewis-resonance-vsepr:lewis}{pasangan bebas} menuju inti. Pada \ce{NH3},
nitrogen adalah ujung negatif setiap ikatan: dipol-dipol ikatan mengarah
dari N ke hidrogen-hidrogen, searah dengan sumbangan \omterm{def:b1:lewis-resonance-vsepr:lewis}{pasangan bebas}, dan
saling menjumlah. Pada \ce{NF3}, fluorin adalah ujung negatif: dipol-dipol
ikatan mengarah dari fluorin-fluorin ke N, berlawanan dengan sumbangan
\omterm{def:b1:lewis-resonance-vsepr:lewis}{pasangan bebas}, dan keduanya hampir saling meniadakan.
\end{solution}

\begin{solution}{exo:b1:lewis-resonance-vsepr:11}
$\mu_b = \mu/(2\cos(\theta/2)) = 1.63/(2\cos 59.75^\circ) =
\qty{1.62}{\debye} = \qty{5.40e-30}{C.m}$. Maka $\delta = \mu_b/(e\,d) =
5.40 \times 10^{-30}/(1.602 \times 10^{-19} \times 143.2 \times 10^{-12})
= 0.24$.
\end{solution}

\begin{solution}{exo:b1:lewis-resonance-vsepr:12}
\ce{N=N=O} dengan \omterm{def:b1:lewis-resonance-vsepr:formal-charge}{muatan formal} $-1$, $+1$, $0$; \ce{N#N-O} dengan $0$,
$+1$, $-1$. Ikatan \ce{N-N} (\qty{112.8}{pm}) mendekati ikatan rangkap
tiga \ce{N2} (\qty{109.8}{pm}): struktur \ce{N#N-O} lebih berbobot, dan
struktur itu menempatkan muatan negatif pada oksigen, atom yang lebih
elektronegatif. Namun, ikatan \ce{N-O} lebih pendek daripada ikatan
tunggal: struktur yang lain juga turut menyumbang.
\end{solution}

\begin{solution}{pb:b1:lewis-resonance-vsepr:1}
\textbf{1.} \ce{N2} 10, \ce{NO} 11, \ce{NO2} 17, \ce{N2O} 16, \ce{O3} 18,
\ce{HNO3} 24.
\textbf{2.} \ce{N#N} dengan satu \omterm{def:b1:lewis-resonance-vsepr:lewis}{pasangan elektron bebas} pada setiap
nitrogen.
\textbf{3.} \ce{N=O} dengan dua \omterm{def:b1:lewis-resonance-vsepr:lewis}{pasangan bebas} pada oksigen, satu
\omterm{def:b1:lewis-resonance-vsepr:lewis}{pasangan bebas} dan satu elektron tunggal pada nitrogen. Dengan 11
elektron, bilangan ganjil, selalu ada satu atom yang jumlah elektronnya
ganjil: nitrogen dikelilingi 7 elektron.
\textbf{4.} \ce{N=N=O}: $-1$, $+1$, $0$; \ce{N#N-O}: $0$, $+1$, $-1$.
\textbf{5.} \ce{O=O-O}: oksigen pusat $+1$ (satu \omterm{def:b1:lewis-resonance-vsepr:lewis}{pasangan bebas}, tiga
ikatan), oksigen ujung yang berikatan tunggal $-1$, yang lain 0.
\textbf{6.} Seperti pada \cref{ex:b1:lewis-resonance-vsepr:hno3}:
nitrogen $+1$, oksigen yang berikatan tunggal dan tanpa hidrogen $-1$,
yang lain 0.
\textbf{7.} \ce{NO} dan \ce{NO2} (11 dan 17 elektron).
\textbf{8.} Dua struktur ekuivalen, dengan ikatan rangkap dua pada salah
satu sisi: orde 1.5 untuk setiap ikatan.
\textbf{9.} Ya: \qty{127.8}{pm} terletak di antara ikatan rangkap dua
(\qty{120.8}{pm}) dan ikatan tunggal (\qty{147.5}{pm}).
\textbf{10.} Ikatan rangkap dua bergiliran pada ketiga oksigen; setiap
ikatan \ce{N-O} rangkap dua pada satu dari tiga struktur: orde $4/3$.
\textbf{11.} Ikatan yang panjang, \qty{140.6}{pm}, adalah \ce{N-OH},
ikatan tunggal. Kedua oksigen tanpa hidrogen berbagi satu ikatan $\pi$
melalui dua \omterm{def:b1:lewis-resonance-vsepr:resonance}{struktur resonansi} yang ekuivalen: dua ikatan berorde 1.5,
121.1 dan \qty{119.9}{pm}.
\textbf{12.} Nitrogen membawa satu elektron tunggal, bukan \omterm{def:b1:lewis-resonance-vsepr:lewis}{sepasang
elektron bebas}: elektron itu menolak pasangan-pasangan ikatan lebih
lemah daripada sepasang elektron, dan ikatan-ikatan itu membuka melewati
$120^\circ$.
\textbf{13.} \ce{N#N-O}, dengan muatan negatif pada oksigen.
\textbf{14.} \ce{O3}: \ce{AX2E}, bengkok. \ce{N2O}: \ce{AX2} (N pusat),
linear. \ce{NO3-}: \ce{AX3}, segitiga datar. Nitrogen \ce{HNO3}:
\ce{AX3}, segitiga datar.
\textbf{15.} \omterm{def:b1:lewis-resonance-vsepr:lewis}{Pasangan elektron bebas} oksigen pusat menolak pasangan-pasangan
ikatan lebih kuat daripada tolakan di antara pasangan-pasangan ikatan
itu sendiri.
\textbf{16.} Ozon bengkok dan atom pusatnya berbeda dari atom-atom ujung
(\omterm{def:b1:lewis-resonance-vsepr:formal-charge}{muatan formal} positif di tengah, muatan negatif terbagi pada kedua
ujung): sebuah dipol kecil sepanjang garis bagi sudutnya. \ce{N2O}
linear, tetapi kedua ujungnya atom yang berbeda: sebuah dipol kecil.
\ce{N2} simetris: tanpa dipol.
\textbf{17.} \ce{CO2} linear (\ce{AX2}): dipol-dipol ikatannya saling
meniadakan. \ce{SO2} bengkok (\ce{AX2E}): dipol-dipolnya saling
menjumlah.
\textbf{18.} \ce{AX3E}: \omterm{def:b1:lewis-resonance-vsepr:lewis}{pasangan bebas} menyempitkan sudut-sudutnya di
bawah $109.5^\circ$; terukur $106.7^\circ$.
\textbf{19.} \ce{AX2E2}: dua \omterm{def:b1:lewis-resonance-vsepr:lewis}{pasangan bebas} menyempitkannya lebih jauh;
$104.5^\circ$.
\textbf{20.} $\mu = 2\mu_{OH}\cos(\theta/2)$
(\cref{prop:b1:lewis-resonance-vsepr:dipole-sum}).
\textbf{21.} $\mu_{OH} = 1.857/(2\cos 52.24^\circ) = 1.857/1.2246 =
\qty{1.52}{\debye}$.
\textbf{22.} $1.516 \times 3.336 \times 10^{-30} = \qty{5.06e-30}{C.m}$.
\textbf{23.} $e\,d = 1.602 \times 10^{-19} \times 95.8 \times 10^{-12} =
\qty{1.535e-29}{C.m}$, yaitu \qty{4.60}{\debye}.
\textbf{24.} $\delta = 5.06 \times 10^{-30}/1.535 \times 10^{-29} =
\textbf{0.33}$: ikatan \ce{O-H} air membawa sepertiga muatan yang akan
dibawa oleh ikatan yang sepenuhnya ionik.
\end{solution}
